The critical, higher-speed signals, like DCMI, clocks, memory and SDMMC were all routed first. This helped prioritise signal integrity on these critical traces since I could utilise the entire PCB. Moreover, this helped with the placement of components, since I was able to identify early on where free space exists and where I can potentially route traces. The DCMI signals were length matched to within $\pm5$mm and QSPI was delay matched within 50ps (DCMI did not require length matching at this speed according to \href{https://www.reddit.com/r/AskElectronics/comments/1o9rrs0/is_there_a_required_delaylength_tolerance_for/}{the community}, but it was included just as a precaution). Moreover, \href{https://www.reddit.com/r/PCB/comments/1p12ws5/recommended_length_matching_properties_altium/}{the community} also suggested that Altium's default accordion/sawtooth pattern for length matching is okay for the current frequency application ($\leq160$MHz).
\nl
90° bends were avoided at all costs to prioritise signal integrity and avoice signal reflections, instead, two 45° turns were used. Where applicable, traces were set to exit and enter pads at 90° to avoid any potential manufacturing issues. Teardrops were used to ease transitions between thin traces and larger pads. This was accomplished using Altium Designer's built-in teardrop generator, where the teardrop style was set to 'curved' and only 'Adjust teardrop size' was selected.
\nl
Where applicable, I avoided layer transitions to prioritise signal integrity. This helped keep as many traces as possible on the same layer, and avoid needing transfer vias which take up extra space. Lastly, the analog signals were kept away from digital signals to minimise the impact of coupling.

\paragraph{Remark}
Delay matching the QSPI and DCMI traces were not necessary (as mentioned), but were harmless if implemented correctly. I had not adhered to proper accordion design rules when greating the traces, so there was a high likelihood that the signal integrity was slightly worse. As the signal speeds were relatively low (less than 500 MHz), it's likely that this bad design did not impact the signals greatly, but it is greatly advised if delay matching is used, that it is implemented correctly. This design aspect fell through the cracks, even during PCB inspection by Tristan Davies.